At the baseline, every ablation reduces the broker's ranking advantage.  The
size-matched gap is \rav{size_matchedBaseline_rank_gap}, a paired change of
\rav{size_matchedBaselineDelta_rank_gap} from paired Ridge (95\% Monte Carlo
interval [\rav{size_matchedBaselineDeltaLower_rank_gap},
\rav{size_matchedBaselineDeltaUpper_rank_gap}]).  The single-principal and
additive gaps are \rav{single_principalBaseline_rank_gap} and
\rav{additiveBaseline_rank_gap}, with paired changes of
\rav{single_principalBaselineDelta_rank_gap} and
\rav{additiveBaselineDelta_rank_gap}.  Their 95\% Monte Carlo intervals are
[\rav{single_principalBaselineDeltaLower_rank_gap},
\rav{single_principalBaselineDeltaUpper_rank_gap}] and
[\rav{additiveBaselineDeltaLower_rank_gap},
\rav{additiveBaselineDeltaUpper_rank_gap}].

The size-matched contrast is consistent across the design.  Its rank gap is
lower than paired Ridge in all \rav{nConditions} effective realizations.  The
median change is \rav{size_matchedDeltaMedian_rank_gap}, and its 10th and 90th
percentiles are \rav{size_matchedDeltaQ10_rank_gap} and
\rav{size_matchedDeltaQ90_rank_gap}.  The gap nevertheless remains positive in
\rav{size_matchedPositiveN_rank_gap}/\rav{nConditions} realizations.  The gap
change is not solely a change in the broker's prediction.  The median broker
rank correlation changes by \rav{size_matchedDeltaMedian_broker_holdout_rank},
while the median agent rank correlation changes by
\rav{size_matchedDeltaMedian_agent_holdout_rank}.  The agent rank correlation
also exceeds paired Ridge in all \rav{size_matchedGreaterN_agent_holdout_rank}
realizations under size matching, all
\rav{single_principalGreaterN_agent_holdout_rank} under single-principal
history, and \rav{additiveGreaterN_agent_holdout_rank} under additive Ridge.
This pattern is consistent with an exploration mechanism: altered broker
rankings change the partners agents encounter and may broaden the diversity of
partner types in their learning histories.  The saved diagnostics do not
measure that diversity, so these results do not establish the proposed
mechanism directly.

The pair-versus-additive contrast holds both endpoint types in each broker
fitting row while removing the bilinear pair features.  As
Figure~\ref{fig:ridge-ablation-grid} shows, the system-level effect is larger
when the interaction component of the production function is strong.  The
median additive-Ridge gap rises from
\rav{additiveRho0_rank_gap} at $\rho=0$ to
\rav{additiveRho0p85_rank_gap} at $\rho=0.85$ and
\rav{additiveRho1_rank_gap} at $\rho=1$.  At $\rho=1$, where $\delta$ drops out
and the additive representation is correctly aligned with the production
function, additive Ridge exceeds the paired-Ridge gap of
\rav{pairRho1_rank_gap}.  The latter result is consistent with a finite-sample
cost from unnecessary bilinear regressors, but the rank-gap contrast alone
does not establish that mechanism because matching and agent learning also
adjust.

Single-principal and additive Ridge both produce additive pair scores, but they
use different fitting rows.  Additive Ridge retains $x_i+x_j$ in each row;
single-principal Ridge retains one randomly selected endpoint and reconstructs
a pair score by adding two fitted one-type scores.  Their baseline paired
rank-gap difference (single-principal minus additive) is
\rav{singleMinusAdditiveBaseline_rank_gap}, with a 95\% Monte Carlo interval of
[\rav{singleMinusAdditiveBaselineLower_rank_gap},
\rav{singleMinusAdditiveBaselineUpper_rank_gap}].  Across effective realizations,
the difference has median \rav{singleMinusAdditiveMedian_rank_gap} and 10th and
90th percentiles \rav{singleMinusAdditiveQ10_rank_gap} and
\rav{singleMinusAdditiveQ90_rank_gap}; it is positive in only
\rav{singleMinusAdditiveGreaterN_rank_gap}/\rav{nConditions} realizations.
Thus, among the two additive-score variants, retaining both endpoint types in
each fitting row produces a modest system-level ranking benefit in most
realizations.  This is not a comparison of bilinear representations: neither
variant can represent the bilinear interaction, which is strongest at low
$\rho$.

The observed brokered-minus-self output gap remains positive in all
\rav{nConditions} realizations under each ablation, even when the broker's rank
gap is negative.  This channel mean difference is descriptive, not a causal
estimate of brokerage's effect on output.  Agents endogenously choose channels,
the channels expose them to different candidate sets, and acceptance decisions
differ.  The positive gap therefore does not show that brokerage improves a
match while holding the demander and feasible candidates fixed, nor does it
identify access, selection, or network position as its cause.  The ablation
contrasts are full-system interventions: prediction, matching, learning
histories, outsourcing, and network position adjust together.  They estimate
the consequences of each modeled restriction for the resulting system, not
additive shares of a fixed-data prediction advantage.
